\chapter{Literature Review}

This thesis draws on research in road-event detection, long-tailed recognition, video representation learning, and vision--language classification. These areas address different parts of the problem: detecting an agent, representing its behavior, learning from unequal class frequencies, and sharing information among related event descriptions. The review distinguishes the evidence established in each setting from the questions that remain for rare ROAD-Waymo compositions. In particular, improvements in image classification, video recognition, or unseen attribute--object recognition are motivations for the proposed study, rather than direct evidence of improved road-event detection.

\section{Road-Event Benchmarks and Visual Representations}

Pedestrian-behavior research illustrates why detection alone is insufficient for scene understanding. JAAD pairs pedestrian boxes with behavioral and contextual annotations, enabling study of crossing behavior and interaction with traffic \cite{jaad-crossing}. PIE adds human-referenced crossing-intention estimation and future trajectory prediction \cite{pie}. These tasks connect appearance, motion, and context, but their targets must be distinguished: an observed crossing action, a future crossing outcome, and a human judgment of intention are different labels. ROAD extends the event-recognition setting across road-agent categories through agent, action, and location compositions. This thesis adopts that event vocabulary and does not transfer pedestrian-intention interpretations to ROAD-Waymo action labels.

Driving datasets differ in the relationships their annotations support. The Waymo Open Dataset provides synchronized camera and LiDAR observations with object annotations for perception research across varied locations \cite{waymo}. BDD100K supports heterogeneous driving tasks, including detection and other forms of scene understanding \cite{bdd}. These benchmarks establish the importance of localization and diverse observations, but object identity alone does not describe the event in which an agent participates. AVA offers a complementary action-detection setting: it associates spatially and temporally localized people with multiple atomic action labels \cite{ava}. Its relevance is the need to classify an actor within a scene rather than assign a single label to an entire video. Its movie-based action vocabulary and evaluation setting differ from those of road-event detection.

ROAD explicitly represents road events through agents, actions, locations, and their compositions. Singh et al. introduced the benchmark and a 3D-RetinaNet baseline for event detection \cite{road}. Duplex labels combine agent and action, while triplets also include location. ROAD-Waymo applies compatible event annotations to US driving video from Waymo Open, allowing comparison with the original UK ROAD data within the ROAD++ framework \cite{road-waymo}. The present thesis uses ROAD-Waymo rather than combining the two domains. Consequently, evidence about geographical transfer and evidence about rare compositions within ROAD-Waymo must remain separate.

ROAD-R adds explicit logical requirements to ROAD and studies learning and prediction mechanisms that use those requirements \cite{road-r}. This distinguishes two meanings of structure. Logical requirements specify which combinations are admissible; semantic representations describe similarities between class meanings. Encouraging constraint satisfaction during training is also different from enforcing admissibility at prediction time. A classifier may produce logically compatible outputs while failing to recognize the correct event. Conversely, improved average precision does not guarantee consistency between primitive and composition predictions. The proposed semantic classifier therefore addresses recognition and does not claim the guarantees of a constrained-output method.

The replicated detector connects several foundational visual architectures. Residual networks make deep representations easier to optimize through residual mappings \cite{resnet}. Inflated 3D networks extend image filters and pooling operations into the temporal dimension to learn video features \cite{i3d}. Feature Pyramid Networks combine a top-down pathway with lateral connections to obtain representations at several spatial scales \cite{fpn}. RetinaNet couples dense detection with focal loss, which reduces the training contribution of easy examples \cite{focal-loss}. ROAD's 3D-RetinaNet adapts this detection approach to road events \cite{road}. The local baseline uses a ResNet50-I3D backbone, so the I3D citation establishes the inflation principle without implying that the local model is the original paper's exact backbone.

Region features provide a basis for separating localization from classification. Mask R-CNN introduces RoIAlign to preserve spatial alignment when extracting a fixed-size representation for a region \cite{mask-rcnn}. In the developed pipeline, a detector supplies the box and RoIAlign reads a frozen video feature map at that box. This separates whether an object was proposed from what features are available to classify it. It also clarifies why copying scores from an overlapping detection is a different operation: score transfer remains dependent on the original detector producing a suitable match.

Video representations provide further alternatives to a fixed backbone. SlowFast networks use a slow pathway for spatial semantics and a faster pathway for fine temporal information \cite{slowfast}. Video Swin Transformer introduces local spatiotemporal attention through a hierarchical transformer design \cite{video-swin}. VideoMAE learns video representations by reconstructing heavily masked video content, emphasizing self-supervised pretraining \cite{videomae}. These approaches demonstrate that temporal sampling, architecture, and pretraining objectives can all influence the representation. They motivate controlling the visual input when evaluating semantic classifiers: replacing both an encoder and a classifier cannot isolate the effect of language-derived class structure.

\section{Learning Under Long-Tailed Class Distributions}

Class imbalance can affect both the representation and the decision rule learned from it. Focal loss addresses the dominance of easy examples in dense detection \cite{focal-loss}; this is related to, but not identical to, unequal support among foreground event categories. A rare triplet may remain difficult even after background examples receive less influence. Thus, the use of focal loss in the developed pipeline does not by itself establish that rare compositions have been adequately addressed.

Several approaches adjust how class frequency enters the learning objective. Cui et al. derive class-balanced weights from the effective number of samples, reflecting diminishing additional information as the number of examples grows \cite{class-balanced}. Cao et al. introduce a label-distribution-aware margin and defer reweighting until after an initial representation-learning period \cite{ldam}. Both methods use class frequency, but they intervene differently: one changes the contribution of examples, while the other changes class-dependent margins and the training schedule. For the proposed study, these distinctions reinforce the need to define tail membership from training data and to keep that definition fixed during evaluation.

Balanced Meta-Softmax addresses distribution shift between training and evaluation through Balanced Softmax and a complementary sampling mechanism \cite{balanced-softmax}. Menon et al. study frequency-based logit adjustment either after training or within the loss \cite{logit-adjust}. These methods show why a change in class scores may improve a class-balanced metric without changing the underlying features. Their softmax formulations also assume a multiclass setting that differs from the independent sigmoid outputs used for multi-label road events. Applying their conclusions directly to the proposed classifier would require an appropriate multi-label formulation and a separate experiment.

Decoupling offers an experimental way to study these effects. Kang et al. separate representation learning from classifier learning and show that adjusting a classifier over a naturally sampled representation can be effective on long-tailed image benchmarks \cite{decoupling}. BBN uses separate branches and a cumulative learning schedule to balance representation and classifier learning \cite{bbn}. These studies motivate treating the classifier as an explicit component rather than attributing every change to the encoder. The developed road-event pipeline follows this principle through frozen features and separately trained heads, while recognizing that image-level classification and box-level event detection have different supervision and error sources.

Other methods change representation geometry or model diversity. RIDE uses multiple distribution-aware experts to address bias and variance under long-tailed sampling \cite{ride}. Parametric Contrastive Learning introduces learnable class centers to counter the tendency of supervised contrastive learning to favor frequent classes \cite{paco}. These approaches show that shared structure can be imposed through more than language. A semantic classifier should therefore be assessed against a meaningful classifier control, rather than compared only with a weak independent-label baseline. They also motivate reporting common-class behavior, since improvements for rare categories may be accompanied by regressions elsewhere.

VL-LTR is especially relevant because it introduces class-level text descriptions into long-tailed visual recognition \cite{vl-ltr}. It learns visual and linguistic representations and reports benefits for classes with fewer image examples. This provides evidence that language can contribute to recognition under imbalance, but it does not establish the same outcome for multi-label road-event boxes. The proposed study differs in using a fixed benchmark vocabulary, frozen crop features, and fixed phrase embeddings with a supervised projection. Its evaluation must establish whether those design choices help rare triplets under the road-event protocol.

Taken together, these methods identify several possible explanations for a tail improvement: better visual features, changes in sampling or loss, classifier calibration, additional capacity, and information shared among classes. The thesis focuses on learned composition and semantic classifier directions. Frequency-based methods remain relevant alternatives, while the primary ablations are designed to distinguish semantic structure from changes to the representation and classifier.

\section{Vision--Language Adaptation for Video Recognition}

CLIP learns aligned image and text representations through contrastive pretraining, allowing class descriptions to define classifier directions \cite{clip}. The resulting relationship between language and visual content is the basis for considering event phrases as more informative than unrelated class identifiers. However, a description that is linguistically plausible may not distinguish visually similar actions or locations. The relevant question is whether the learned alignment retains the distinctions needed for road-event recognition.

Prompt-learning methods demonstrate both the usefulness and sensitivity of text-based classification. CoOp learns context vectors while keeping the pretrained model fixed \cite{coop}. CoCoOp conditions part of the prompt on the input to improve generalization beyond the classes used for adaptation \cite{cocoop}. MaPLe couples prompts in the visual and language branches to adapt their representations together \cite{maple}. These methods make different parts of the model trainable. They therefore cannot be treated as equivalent to a classifier whose phrase matrix remains frozen. For this proposal, their main implication is that the choice and adaptation of class descriptions must be documented and controlled.

Adaptation can also use labeled visual examples directly. Tip-Adapter builds a key--value cache from a few-shot training set and combines feature retrieval with the pretrained model's predictions \cite{tip-adapter}. It provides a useful contrast to prompt learning: downstream performance may improve through visual evidence without changing the textual description. This motivates keeping the training population and feature source fixed when comparing the proposed semantic head with a free classifier. A difference between those heads can then be interpreted more narrowly than a comparison in which the representations and examples also change.

Transferring image--text models to video introduces temporal design choices. X-CLIP uses cross-frame attention and video-dependent prompting to adapt pretrained image--language representations to video recognition \cite{xclip}. Ju et al. combine continuous prompts with lightweight temporal transformers over frame features \cite{efficient-prompt}. ViFi-CLIP shows that fine-tuning and pooling frame-level CLIP features can provide a strong video baseline without requiring a complex new temporal module \cite{vifi}. Together, these studies caution against attributing an improvement to semantic reasoning merely because a vision--language model is involved. Temporal processing, adaptation, and pooling can each contribute.

InternVideo2 develops a family of video models through progressive training that combines masked video modeling, cross-modal contrastive learning, and generative objectives \cite{internvideo2}. The proposed implementation uses the released InternVideo2\_CLIP\_S model, whose repository supplies its configuration and weights \cite{internvideo2-clip-s}. The family-level paper motivates the representation, while the model release identifies the specific component used in this thesis. The proposed head reads an internal crop descriptor and learns a projection into the phrase space; it does not assume that the internal descriptor is already the released model's final video--text embedding.

The distinction between pretraining and the proposed experiment is central. Contrastive video--text learning has already shaped the frozen encoders. The planned ROAD-Waymo experiment instead trains a supervised classifier using those representations. It tests whether retaining semantic directions is useful for rare benchmark compositions. It does not claim that the supervised head itself performs new contrastive pretraining, or that the presence of text automatically provides temporal reasoning.

\section{Compositional Structure and the Research Gap}

Compositional recognition asks how predictions for parts relate to predictions for their combinations. ROAD considers both direct composite prediction and composition of primitive scores \cite{road}. The developed MLP follows the former goal through a structured input: it receives primitive predictions together with visual features and learns scores for the benchmark's valid duplexes and triplets. This gives it access to relationships that a fixed product or minimum of primitive scores cannot adapt to the task. It does not impose logical agreement between independently predicted heads.

Compositional zero-shot learning provides related evidence about sharing information among concepts. Naeem et al. connect states, objects, and their compositions through a graph so that relationships among primitives inform composition classifiers \cite{cge}. Nayak et al. learn compositional soft prompts in which attribute and object tokens can be recombined to describe unseen combinations \cite{csp}. These methods explicitly target generalization to compositions absent during training. The proposed ROAD-Waymo study instead evaluates a fixed vocabulary that includes rare but observed classes. Their results motivate structured sharing, but do not justify a zero-shot claim for this thesis.

The literature therefore supports a focused gap rather than the claim that semantic classification is new in general. Long-tailed recognition studies frequency-sensitive losses, decoupled classifiers, and linguistic information; video adaptation studies the transfer of image--text representations; compositional learning studies shared primitives and combinations. What remains to be established here is whether fixed phrase-based classifier directions improve rare ROAD-Waymo triplets when the visual inputs, candidate boxes, training examples, and evaluation population are controlled.

The developed pipeline supplies the comparison platform. The proposed experiments will relate tail AP to overall and common-class AP, inspect per-class changes, and repeat classifier training with recorded seeds. A positive result must be supported by the relevant ablations, while a negative result can identify limitations in the visual evidence or its alignment with event descriptions. This framing connects the related work to an empirical question that the proposed experiment can answer without assuming that semantic structure will necessarily help.
